

\chapter{Introduction}

 Let us try to imagine that someone has to get from IIT Madras to
the airport, and they have $4$ real choices like Car, Bike, taxi, bus..
And so, based on the cost, comfort, speed they have some natural
preference ranking among this $4$, let's say they choose or pick
\begin{equation*}
    \text{Car} > \text{Taxi} > \text{Bike} > \text{Bus}
\end{equation*}

Now, suppose we add a $5^{\text{th}}$, a clear irrelevant option to the
list which would almost take $3$ hours and nobody would seriously
choose this. Common sense tells us that adding `walking' doesn't
really change the feel about the car v/s the taxi v/s the bike, as it
is the irrelevant alternative, and their \textbf{relative preference} between all these $4$ alternative would not change at all.

So this idea is called independence of irrelevant alternatives and the
core behind the classical model of choice.

But the core idea is that maybe human brains may not be registering
the preference in the same way of that of multinomial logits. As
there is a strong experimental evidence proves that it would shift the attractiveness of things in the way that how attractive the existing options really looks $\rightarrow$ this is actually called violations of independence of irrelevant alternatives.\cite{seshadri2019}